\chapter{Analisis sobre parametros de hoja de datos}

\section{Características Eléctricas (Regímenes máximos)}
\begin{itemize}
  \item \textbf{VRRM} (Repetitive Peak Reverse Voltage): Máxima tensión pico inversa repetitiva que el diodo puede soportar sin dañarse durante su operación normal.
  \item \textbf{VRSM} (Non-Repetitive Peak Reverse Voltage): Máxima tensión pico inversa no repetitiva que el diodo puede soportar (por ejemplo, durante transitorios o picos de voltaje excepcionales).
  \item \textbf{IFRMS} (Forward RMS Current): Corriente eficaz (RMS) máxima que el diodo puede conducir en polarización directa de forma continua sin exceder sus límites térmicos.
  \item \textbf{IFSM} (Surge Forward Current): Corriente pico máxima que el diodo puede soportar en forma de pulsos cortos, como por ejemplo un pico de corriente de arranque o una sobretensión momentánea.
\end{itemize}

\section{Características Principales}
\begin{itemize}
  \item \textbf{VZ} (Zener Voltage): Tensión de zener o tensión de avalancha en diodos zener, la tensión a la cual el diodo comienza a conducir en sentido inverso manteniendo un voltaje casi constante.
  \item \textbf{PT} (Total Power Dissipation): Potencia máxima que el dispositivo puede disipar sin que su temperatura aumente más allá del límite permitido.
  \item \textbf{TJ} (Junction Temperature): Temperatura máxima permitida en la unión del dispositivo para su correcto funcionamiento y seguridad.
  \item \textbf{VF} (Forward Voltage): Caída de tensión en el diodo cuando está en polarización directa y conduce corriente.
  \item \textbf{IZ} (Zener Current): Corriente que circula por un diodo zener cuando está en la región de ruptura, manteniendo el voltaje de zener constante.
\end{itemize}

\section{Regímenes característicos}
\begin{itemize}
  \item \textbf{VBR} (Breakdown Voltage): Tensión de ruptura, el voltaje inverso a partir del cual el diodo comienza a conducir una corriente significativa en polarización inversa.
  \item \textbf{IR} (Reverse Current o Leakage Current): Corriente de fuga en polarización inversa, corriente muy pequeña que circula cuando el diodo está bloqueando.
  \item \textbf{VF} (Forward Voltage): Igual que antes, la caída de tensión en conducción directa.
  \item \textbf{IRM} (Maximum Reverse Current): Corriente inversa máxima permitida bajo condiciones de prueba específicas sin dañar el dispositivo.
  \item \textbf{IF} (Forward Current): Corriente directa de operación continua del diodo.
\end{itemize}

\section{Características Térmicas}
\begin{itemize}
  \item \textbf{TJ} (Junction Temperature): Temperatura máxima de la unión.
  \item \textbf{TSTG} (Storage Temperature): Rango de temperatura para almacenamiento seguro del dispositivo sin daño.
  \item \textbf{TA} (Ambient Temperature): Temperatura ambiente máxima de operación del dispositivo.
  \item \textbf{Rthjc} (Thermal Resistance Junction-to-Case): Resistencia térmica entre la unión y la carcasa, mide la eficiencia con la que el calor generado en la unión es conducido hacia el exterior.
\end{itemize}

\section{Características de uso}
\begin{itemize}
  \item \textbf{VRMS} (Maximum RMS Voltage): Valor máximo de tensión RMS que el dispositivo puede soportar (típico en diodos de potencia o rectificadores).
  \item \textbf{Ptot} (Total Power): Potencia total máxima que puede disipar el dispositivo.
  \item \textbf{Trr} (Reverse Recovery Time): Tiempo que tarda el diodo en dejar de conducir después de pasar de polarización directa a inversa, importante en aplicaciones de conmutación rápida.
\end{itemize}
